% 立直（リーチ）：聴牌（テンパイ）を宣言すること

\documentclass[../nankiru_main.tex]{subfiles}

% override main tex render option
\renewcommand{\renderOption}{1}

\begin{document}

\subfilepreamble{立直判断}

\subsection{先制両面立直}

\begin{table*}[t]
  \centering
  \begin{tabular}{l|SSSSSS}
    巡目       & \multicolumn{2}{c}{− 5巡目 −} & \multicolumn{2}{c}{− 8巡目 −} & \multicolumn{2}{c}{− 12巡目 −} \\
    手牌基本価値 & 立直 & ダマ & 立直 & ダマ & 立直 & ダマ \\
    \hline
    のみ手     & 1300 &      & 600 &     & -100 &  \\
    1飜30符    & 1900 & 600 & 1100 & 200 & 300 & -300 \\
    1飜40符    & 2800 & 1000 & 1900 & 500 & 900 & 0 \\
    2飜30符    & 3600 & 1600 & 2600 & 1000 & 1400 & 400 \\
    2飜40符    & 4500 & 2300 & 3400 & 1600 & 1900 & 800 \\
    3飜30符    & 5500 & 3400 & 4200 & 2600 & 2600 & 1600 \\
    3飜40符    & 5600 & 4800 & 4300 & 3800 & 2600 & 2500 \\
    3飜50符    & 5600 & 5600 & \color{red}4300 & \color{red}4500 & \color{red}2600 & \color{red}3000 \\
    4飜30符    & 7000 & 6300 & 5500 & 5100 & \color{red}3500 & \color{red}3600 \\
    満貫確定    & 7000 & 6500 & 5500 & 5300 & \color{red}3500 & \color{red}3700 \\
    立直で跳満 & 8300 & 7500 & 6500 & 6100 & \color{red}4300 & \color{red}4400 \\
    跳満確定    & \color{red}9600 & \color{red}9900 & \color{red}7600 & \color{red}8200 & \color{red}5100 & \color{red}5900
  \end{tabular}
  \caption{先制両面立直とダマの局収支\citep[][p. 23]{miinin-stat}。赤字はダマの方が有利の場面}
  \label{tab:ev-ryanmen-riichi}
\end{table*}

表~\ref{tab:ev-ryanmen-riichi} は先制立直とダマの局収支の比較。例え、手牌に1飜40符があって、5巡目で立直するとその局収支は2800点、一方ダマにするとその局収支は1000点しかないので、早めに先制立直を打つのは大体正解。中盤（8巡目ごろ）では、手牌の価値は\ref{mj:rch-1-1}の3飜50符\footnote{門前加符30符 ＋么九牌暗刻8符＋么九牌暗刻8符＝50符（切り上げ）。詳細な点数計算は\ref{sec:point}参照。}より少ないなら立直するのは大体正解。

\tehai{12388m 67p 555 666z}[東1局　南家　8巡目　ドラ \mj{1m}][rch-1-1]

終盤（12巡目以降）で\ref{mj:rch-1-2}のような3飜40符手牌があるなら、立直とダマの局収支に大きな差は出ないから、アガリ率を優先してダマにするのがおすすめ。

\tehai{233445m 2267p 888s}[東1局　南家　12巡目　ドラ \mj{2p}][rch-1-2]

\subsubsection*{3枚見えピンフ} 

\textbf{待ちが3枚見えてても立直する}\citep[][pp. 93--94]{miinin-choice}。手牌\ref{mj:rch-1-3}のように待ちが真ん中、\mj{36p} が3枚見えていても\footnote{\mjfn{7p} がドラなら \mjfn{6p} がドラ指示牌として1枚使っている。}、立直の局収支がダマより高い。

\tehai{345678m 3345p 123s}[東1局　南家　8巡目　ドラ \mj{7p}][rch-1-3]


\subsubsection*{親ドラポンに対する}

\textbf{親ドラポンが入ってもドラ2両面で勝負する}\citep[][pp.~136--137]{miinin-choice}。ベタオリなら局収支が$-$1900になるので価値があれば11巡目でも両面立直で勝負する方は局収支（$-$1100）が高い。よって、手牌\ref{mj:rch-1-4}なら親のドラ \mj{7z} ポンを無視して立直すべき。ただし、ドラが1枚だけならベタオリすべき。

\tehai{307m 789p 2206s 333z -6m}[東1局　南家　11巡目　ドラ \mj{7z}][rch-1-4]

\subsubsection*{両面リーのみ}

\tehai{34m 234p 3335567s 2 -2z}[東1局　西家　3巡目　ドラ \mj{7m}][2-4-0]

\tehai{34m 2345p 333567s 2 -2z}[東1局　西家　3巡目　ドラ \mj{7m}][2-4-1]

\tehai{34m 2346p 333567s 2 -2z}[東1局　西家　3巡目　ドラ \mj{7m}][2-4-2]

全ての状況は早めの3巡目だが、そのままで立直すれば打点はただの1300。手牌\ref{mj:2-4-0}は \mj{22z} を落としタンヤオや三色やピンフで打点は大きく上昇する。手牌\ref{mj:2-4-1}は \mj{22z} を落として \mj{25m} を引けばまだ三色やピンフに替わり機会が多いのでダマはおすすめ。\ref{mj:2-4-2}が \mj{6p} 単騎になる可能性があるので、それを避けて即リーがおすすめ\citep[][pp. 68--73]{tetsusenryaku}。


\subsubsection*{三面張と三色確定両面の二択}

手牌\ref{mj:ssk-25}は打 \mj{6p} により三面張と打 \mj{3p} による456の三色確定。

\tehai{22456m 345667p 456s}[8巡目][ssk-25]
% \citep{clearrain-ssk}
% ドラなし -> 3p
% ドラ2枚 -> 2m


%--------------------------------------------
\subsection{先制愚形立直} \label{subsec:senseigukei}

\begin{table*}[t]
  \centering
  \begin{tabular}{l|SSSSSS}
    巡目       & \multicolumn{2}{c}{− 5巡目 −} & \multicolumn{2}{c}{− 8巡目 −} & \multicolumn{2}{c}{− 12巡目 −} \\
    手牌基本価値 & 立直 & ダマ & 立直 & ダマ & 立直 & ダマ \\
    \hline
    のみ手     & 500 &      & -100 &      & -600 & \\
    1飜40符    & 1600 & 500 & 800 & 0 & 0 & -500 \\
    2飜40符    & 2900 & 1600 & 1900 & 900 & 800 & 100 \\
    3飜40符    & 3800 & 3800 & 2700 & 2600 & 1400 & 1300 \\
    3飜50符    & 3800 & 4500 & 2700 & 3100 & 1400 & 1700 \\
    満貫確定    & 4800 & 5300 & 3500 & 3700 & 1900 & 2100 \\
    立直で跳満 & 5900 & 6100 & 4300 & 4400 & 2600 & 2600 \\
    跳満確定    & 6900 & 8200 & 5200 & 6000 & 2100 & 3700
  \end{tabular}
  \caption{先制愚形立直とダマの局収支\citep[][p. 28]{miinin-stat}}
  \label{tab:ev-gukei-riichi}
\end{table*}

\citet[p. 165]{rb1}によると、親なら、もしくは立直しなくても手牌にもう1飜がついていたら、愚形でも立直すべき。表~\ref{tab:ev-gukei-riichi} によると、巡目に関わらず、手牌の価値は3飜40符（ダマで5200、立直で満貫）があればダマしてもいい。例え\ref{mj:rch-2-1}のような手牌は立直してもダマしても局収支に大きな差がついていないので和了率優先ならダマにしよう。

\tehai{112233m 22789p 68s}[東1局　南家　8巡目　ドラ \mj{2m}][rch-2-1]


\subsubsection*{愚形のみ手} 

\citet[p. 165]{rb1}によると、子でドラなし愚形のみ手を立直してはいけない。\citet[p. 30]{miinin-stat}もそういう手を非推奨する。そうすると選択肢には（ア）聴牌とってダマ（イ）聴牌外し（ウ）ベタオリ、三つある。（ウ）の局収支は-1500点のでそれも除ける\citep[][p. 30]{miinin-stat}。（ア）の局収支は表~\ref{tab:ev-tsumosen} の通り、必ず0以下になる。

\begin{table}[ht]
  \centering
  \begin{tabular}{llll}
    巡目 & 5巡目 & 8巡目 & 12巡目 \\
    \hline
    局収支 & 0 & $-$300 & $-$700
  \end{tabular}
  \caption{役なしドラなし愚形（ツモ専）の局収支\citep[][p. 32]{horiuchi}}
  \label{tab:ev-tsumosen}
\end{table}

よって、通常の愚形よりいい待ち（筋待ち・1289待ち）の場合、\citet[p. 30]{miinin-stat}は「立直をかけて」と主張し、他の愚形（特に456待ち）は「終巡ならばダマして形式聴牌を狙い、中巡より前であれば聴牌外し」の選択を提出した。


\subsubsection*{真ん中の愚形} 

456待ちの愚形は1289待ちよりアガリにくい。けれど、\citet{nisi-456kanchan}のシミュレーションによると中盤に入ったら真ん中の愚形でも立直すればダマより局収支が高い。\citet{ayataka}によると「先制であれば多くの場合でリーチ有利」。\citet[pp.~31--32]{ishibashi-black}によると無筋46リーのみはしない方がいいが、5待ちならツモ赤ドラの場合があるなら立直する。


\subsubsection*{手変わりの有効局面} 

前に触られた「聴牌外し」は、いつも序盤で立直より有利なわけではない。例え\ref{mj:rch-2-2}のようなツモ専の手牌は、\mj{7s} 切りで立直すれば1300点、\mj{1p} 切ればくっつき一向聴になる。しかし、リーのみと聴牌外しとの和了率の差は5巡目ではすでに5％あって、巡目が進むほどこの差は大きくなる\citep[][p. 34]{miinin-stat}。よって、聴牌を外して良形になれそうでも先制立直をするのは正解。

\tehai{234999m 1277p 4567s}[東1局　南家　5巡目　ドラ \mj{6z}][rch-2-2]

聴牌外しにより打点上昇のケースは\ref{mj:rch-2-3}のように辺張を払えばタンヤオがつく場面。中盤でも聴牌外しの方が若干有利\citep[][p. 35]{miinin-stat}。しかし、\mj{2m} がドラになって手牌に既に1飜があるなら、序盤で聴牌外しの方が有利だが、中盤では場況によって判断する。特に\ref{mj:rch-2-4}のような形は別に強くない場合は立直に寄る。

\tehai{234888m 1288p 4567s}[東1局　南家　5巡目　ドラ \mj{6z}][rch-2-3]

\tehai{234888m 1288p 4888s}[東1局　南家　8巡目　ドラ \mj{2m}][rch-2-4]

要注意のは序盤でドラ1の愚形のみ手の場合、手変わりができれば、\citet[p. 35]{miinin-stat}の主張は\citet[p. 165]{rb1}と異なる。前者は「聴牌外し」を推奨し、後者は「即リー」を推奨する。


\subsubsection*{当たり牌の種類と残り枚数} 

\begin{table*}[t]
  \centering
  \begin{tabular}{l|llllllll}
               & \multicolumn{2}{l}{全体} & \multicolumn{2}{l}{場に0枚見え} & \multicolumn{2}{l}{場に1枚見え} & \multicolumn{2}{l}{場に2枚見え} \\
               & 和了 (\%) & ツモ (\%) & 和了 (\%) & ツモ (\%) & 和了 (\%) & ツモ (\%) & 和了 (\%) & ツモ (\%) \\
    \hline
    無筋456 & 34 & 62 & 37 & 63 & 30 & 61 & 23 & 57 \\
    無筋37  & 40 & 51 & 43 & 52 & 36 & 49 & 28 & 46 \\
    無筋28  & 43 & 46 & 46 & 47 & 41 & 46 & 34 & 42 \\
    片筋456 & 36 & 52 & 40 & 53 & 32 & 52 & 24 & 48 \\
    無筋19  & 46 & 36 & -- & -- & 47 & 37 & {\color{red}44} & 33 \\
    筋37    & 48 & 36 & 51 & 37 & 44 & 34 & 35 & 32 \\
    筋28    & 51 & 31 & 55 & 32 & 50 & 31 & {\color{red}43} & 30 \\
    両筋456 & 47 & 33 & 51 & 34 & 42 & 33 & 36 & 31 \\
    筋19    & 57 & 24 & -- & -- & 58 & 24 & {\color{red}56} & 25
  \end{tabular}
  \caption{当たり牌の種類・枚数別5〜12巡目の先制単騎・嵌張立直の和了率とツモの割合\citep[][p. 38]{miinin-stat}。赤字は12巡目の先制両面立直より和了率が高い立直}
  \label{tab:ev-gukei-riichi-machi}
\end{table*}

基準として、12巡目の先制両面立直の和了率は42\%、そのときのツモの割合は54\%\citep[][p. 26]{miinin-stat}。表~\ref{tab:ev-gukei-riichi-machi} によると、自分から2枚見えても、基準より和了率が高いのは無筋19・筋19・筋28。これらの待ちになったら立直をかけるのは正解。

\textbf{残り枚数が1枚減っていても打点が高い方にする}\citep[][pp. 91--92]{miinin-choice}。例え\ref{mj:rch-2-5}には嵌 \mj{6m} と嵌 \mj{4m} の二択がある。そして嵌 \mj{4m} にすれば一盃口がつく。残り4枚の方のアガリ率が高いけど、局収支として嵌 \mj{4m} の方が高い。

\tehai{334557m 222p 789s 33z}[東1局　南家　8巡目　ドラ \mj{9s}][rch-2-5]

残り枚数が2枚しかないなら立直の基準が厳しくなる。けど、場況・ドラの枚数により立直することが得の場面もある\citep{yoteru-twoleft}。表~\ref{tab:gukei-twoleft} の通り、いい形（四連形・中ぶくれなど）がないなら立直の基準が下回って、一方愚形を外しても他の部分で搭子ができるならそこに頼って立直の基準が上回る。また、先制された場合はもっと厳しくなり、ドラが3枚以上持っていなければ安易に立直をかけるべきではない。

\begin{table*}[t]
  \centering
  \begin{tabular}{l|lll}
    ドラ枚数 & いい形がない & いい形がある & 先制された場合 \\
    \hline
    なし & ダマ & 聴牌外し & 追いかけない \\
    1枚 & 待ちがいいなら立直する & 待ちがいいなら立直する & 追いかけない \\
    2枚 & 立直 & 待ちがいいなら立直する & 待ちがいいなら立直する \\
    3枚以上 & 立直 & 立直 & 立直
  \end{tabular}
  \caption{愚形待ちが残り2枚しかない状況の立直判断（七対子除き）}
  \label{tab:gukei-twoleft}
\end{table*}

表~\ref{tab:gukei-twoleft} の中で「待ちがいい」というのは以下の要素で決める、

\begin{itemize}
\item 愚形の色の安さ
\item 山読み
\item 河の筋にかかっている
\end{itemize}

こうすると、ドラなしで\ref{mj:twoleft}は \mj{8p} 落としてマンズの部分を伸ばす方がいい。即リーの場面は \mj{7s} がドラの時。

\tehai{4556m 12377799s 68p}[東1局　南家　8巡目　\mj{7p} が2枚切れ][twoleft]

ちなみに字牌双碰の時2枚切れでも字牌が出やすいのでドラなしでも先制立直すべき。手牌\ref{mj:twoleft-2}は鉄リー。

\tehai{234m 567s 678p 2233z}[東1局　南家　8巡目　\mj{23z} が2枚切れ][twoleft-2]


\subsubsection*{モロ引っ掛け} 

\textbf{筋待ちは無筋よりいい待ち}\citep[][p. 77]{miinin-stat}。待ちが立直宣言牌の筋（モロ引っ掛け）の場合は相手に警戒されやすいが、その和了率はまだ十分。表~\ref{tab:winrate-28} と表~\ref{tab:winrate-37} によると、モロ引っ掛けの和了率は非モロヒより3\%低いが、無筋の方と比べるとそんなに大きな差ではない。よって、\mj{468m} のような赤狙いとモロ引っ掛けと二択があっても \mj{4m} 切って立直するべき\citep[][pp. 71--72]{miinin-choice}。なお、赤を切って筋を作る（赤モロヒ）のは警戒されにくいが、3\%の和了率差に1飜を減らすのは損。

\begin{table}[ht]
  \centering
  \begin{minipage}{.48\linewidth}
    \centering
      \begin{tabular}{ll}
        & (\%) \\
        \hline
        無筋28 & 46 \\
        黒モロヒ28 & 54 \\
        赤モロヒ28 & 57 \\
        非モロヒ & 57
      \end{tabular}
      \caption{嵌28の和了率（5〜12巡目）}
      \label{tab:winrate-28}
    \end{minipage}%
  \begin{minipage}{.48\linewidth}
    \centering
    \begin{tabular}{ll}
      & (\%) \\
      \hline
      無筋37 & 43 \\
      モロヒ37 & 50 \\
      非モロヒ & 53 \\
    \end{tabular}
    \caption{嵌37の和了率（5〜12巡目）}
    \label{tab:winrate-37}
  \end{minipage}%
\end{table}

ただし、「両面＋リャンカン」という良形聴牌の確率は50\%だけなら、\textbf{5を先切って相手を警戒されない}ようにするのは実戦的非効率な打ち方\citep[][pp. 23--25]{ishibashi-black}。例え、手牌\ref{mj:rch-5-1}にタンヤオもピンフもつかないため、良い待ち（良形に限らず）を作るため打 \mj{3m} による50\%の両面聴牌より打 \mj{5p} による100\%のいい待ち（両面＋筋 \mj{2p}）にすべき。

\tehai{233 99m 135p 340 888s}[][rch-5-1]

偶数牌のリャンカン（246や468）の場合なら筋37より中筋5の方が上がりやすい。よって、手牌\ref{mj:rch-5-2}は \mj{2p} 切るべき。

\tehai{233 99m 246p 406s 888s}[8pが河にある][rch-5-2]


\subsubsection*{「双碰不如一嵌」} 

「双碰は嵌張にしかず」とは古来の戦術。嵌張待ちは誰かが3枚持っていてもまだ1枚が残っているから、対死（トイスー）の可能性がある双碰より良い、と考える。だが、同じ無筋で、残り枚数も同じ、手役に絡まない場合、後筋になる確率と待ち種類を考えると\textbf{双碰の方が有利}\citep{tsukemai}。残り枚数が多いの方と筋になっている方を選ぶのは基本。


\subsubsection*{ノベタンとドラ字牌単騎} 

\textbf{ドラ字牌単騎にする}\citep[][pp.~75--76]{miinin-choice}。ノベタンの方がアガリやすいが、ドラ字牌単騎は生牌でも期待値がより高い。


\subsubsection*{ツモり四暗刻}

\textbf{双碰待ちの四暗刻はダマ}\citep[][pp.~179--180]{miinin-choice}。立直してもツモ率が17\%から20\%までしか上がらないので。けれど、点数差がひどい状況なら立直する。


\subsubsection*{山読みによる選択}

\begin{rittai}[h]
  \centering
  
  \drawrittai{
    jicha/seat = 南,
    jicha/hand = 55789m 67p 113789s -8p,
    jicha/river = 1p 3z 4p 5z 9p 3p,
    jicha/riverb = 44z 2m,
    jicha/point = 36400,
    %
    toimen/river = 6z 2s 6z 989p,
    toimen/riverb = 3m 4z,
    toimen/point = 2800,
    %
    kami/hand = XXXXXXX -2'22m111'z,
    kami/river = 8m 211p 7z 3m,
    kami/riverb = 6m 4z 1m 6z 7p,
    kami/point = 46100,
    %
    shimo/river = 9s 1m 3z 8s 7m 8s,
    shimo/riverb = 3s 1m 7z,
    shimo/point = 14700,
    %
    kyoku = 南2局,
    dora = 6p
  }
  
  \caption{Mリーグ2020 1/2 第2回戦 南2局 0本場（ツモ切り暗転なし）\citep[][p.~57]{muhai}}
  \label{rittai:20200102}
\end{rittai}

立体図 \ref{rittai:20200102} のように、\mj{8p} を引いて \mj{5m1s} の双碰待ちと \mj{2s} の単騎待ちを選択する状況である。河によると、下家が2巡前に \mj{3s} を切って、手牌の中に \mj{2s} がいそうである。一方、対面が早めに \mj{2s} を切って、手牌に \mj{1s} がいなさそうである。そして上家が \mj{8m} と \mj{2p} から切って、そのあと真ん中の \mj{3m} と \mj{6m} を切った。もし上家に \mj{1s} があったらそこから切っているであろう\footnote{\mjfn{1s} が二つある可能性なのでは…}。しかも \mj{2m} もポンしたので、ソウズの染め手でもない。これらの情報によると、\mj{5m} はちょっと厳しいけれど \mj{1s} はまだ山にありそうなので、ここで\citet[p.~58]{muhai}は \mj{5m1s} の双碰待ちを選んだ。


%--------------------------------------------
\subsection{愚形の打点が両面より高い}

例として、手牌\ref{mj:rch-6-1}は両面と嵌張の選択。\mj{25m} の両面にして立直をかけるとこの手の価値は2飜30符となり、その局収支は表~\ref{tab:ev-ryanmen-riichi} によると1100点となる。一方、嵌 \mj{2m} にすると三色がついて、表~\ref{tab:ev-gukei-riichi} によると3飜40符の局収支は1900点となる。両面ピンフより嵌張三色の方が高いので、ここで \mj{4m} 切り立直すべき\citep[][p. 55]{miinin-stat}。

\tehai{134m 123p 12356799s}[東1局　南家　8巡目　ドラ \mj{6z}][rch-6-1]

ドラが1枚あるならどうでしょう。両面にすると局収支は2600点となり、嵌張にすると局収支は2700点となる。どっちでもいい場合\citep[][p. 56]{miinin-stat}。

\tehai{134m 123p 12356799s}[東1局　南家　8巡目　ドラ \mj{1p}][rch-6-2]

よって、表~\ref{tab:ev-ryanmen-gukei-compare} の通り、両面待ちには1飜だけなら愚形にして打点を高くにするのは正解。両面待ちにもう3飜がついているなら両面で立直する。両面待ちに2飜があるなら愚形立直と局収支はほぼ同じだが、巡目が深くなるなら愚形ダマの方が有利。

\begin{table*}[t]
  \centering
  \begin{tabular}{l|lllllllll}
    巡目       & \multicolumn{3}{l}{5巡目} & \multicolumn{3}{l}{8巡目} & \multicolumn{3}{l}{12巡目} \\
    手牌価値（両面--愚形）& 両面 & 愚形リー & 愚形ダマ & 両面 & 愚形リー & 愚形ダマ & 両面 & 愚形リー & 愚形ダマ \\
    \hline
    1飜30符--2飜40符    & {\color{red}1900} & {\color{red}2900} & {\color{red}1600} & {\color{red}1100} & {\color{red}1900} & {\color{red}900} & {\color{red}300} & {\color{red}800} & {\color{red}100} \\
    2飜30符--3飜40符    & 3600 & 3800 & 3800 & {\color{blue}2600} & {\color{blue}2700} & {\color{blue}3100} & {\color{blue}1400} & {\color{blue}1400} & {\color{blue}1700} \\
    3飜30符--満貫       & 5500 & 4800 & 5300 & 4200 & 3500 & 3700 & 2600 & 1900 & 2100 \\
  \end{tabular}
  \caption{愚形と両面の比較。特に愚形の打点がより高い場面。赤字は愚形が有利の場面。青字はダマが有利の場面}
  \label{tab:ev-ryanmen-gukei-compare}
\end{table*}


%--------------------------------------------
\subsection{字牌待ち}

\citet[p. 46]{miinin-stat}によると、立直以外1飜以上の価値があれば\textbf{全ての場面では立直が有利}。唯一の要注意は12巡目字牌待ちののみ手の局収支は-300。ダマにすべきの場面は、自分の持ち点が多くて和了率が優先する状況だけ。例え、\ref{mj:rch-4-1}のような跳満確定の手は、自分が四位なら立直して、自分がトップならダマにする。

\tehai{340m 340s 340p 88p 44z}[南3局　南家　8巡目　ドラ \mj{3m}][rch-4-1]

\subsubsection*{両面と字牌待ちの選択} 

\textbf{打点（高目安目があっても）が高いの方を選ぶのは正解}\citep[][p.50]{miinin-stat}。字牌待ちと両面待ちの和了率はほぼ同じだから。例えば、手牌\ref{mj:rch-4-2}を \mj{36p} 待ちにすれば1300点に過ぎない。逆に \mj{4p5z} の双碰にすれば高目が2600点になる。和了率がほぼ同じなら高目があるから局収支も高くなる。しかし、手牌\ref{mj:rch-4-3}のように字牌がオタ風の場合、\mj{36p} の場合にピンフがつくのでそれにするのは正解。

\tehai{234567m 445p 789s 55z}[東1局　南家　8巡目　ドラ \mj{6z}][rch-4-2]

\tehai{234567m 445p 789s 44z}[東1局　南家　8巡目　ドラ \mj{6z}][rch-4-3]

\textbf{どちらの待ちにしても打点が変わらない場合は両面にする}\citep[][p. 52]{miinin-stat}。両面待ちのツモ割合がより高いから、和了時の平均得点もより高い。手牌\ref{mj:rch-4-4}の場合は \mj{2p} 切って立直。

\tehai{234567m 223p 888s 44z}[東1局　南家　8巡目　ドラ \mj{2m}][rch-4-4] 

\subsubsection*{字牌単騎待ちについて枚数別の和了率} 

これまでの話は字牌双碰に関する。七対子などの字牌単騎の場合、\textbf{河に1枚切れの単騎の和了率は最も高い}\citep[][p. 64]{miinin-stat}。表~\ref{tab:winrate-jihai-tanki} によるとその和了率は両面よりも高い。従って、手牌\ref{mj:rch-4-5}は \mj{3z} が1枚切れなら字牌単騎にしよう。でなければ \mj{58p} の両面をとる。

\begin{table}[ht]
  \centering
  \begin{tabular}{llll}
    和了率(\%) & 5巡目 & 8巡目 & 12巡目 \\
    \hline
    字牌単騎（1枚切れ） & 70 & 58 & 45 \\
    両面             & 68 & 57 & 42 \\
    字牌単騎（生牌）   & 65 & 49 & 30 \\
    字牌単騎（2枚切れ）& 55 & 48 & 40
  \end{tabular}
  \caption{枚数別字牌単騎の和了率と両面立直の比較}
  \label{tab:winrate-jihai-tanki}
\end{table}

\tehai{123888m 345s 6788p 3z}[東1局　南家　12巡目　ドラ \mj{1m}][rch-4-5] 

%--------------------------------------------
\subsection{七対子の立直判断}

\textbf{先制なら七対子をダマにするケースはほとんどない}\footnote{最終形が聴牌のデータのみを集計してからの結論なので、手変わりが考慮されていない\citep[][p.~102]{miinin-choice}。ダマにして手変わりによる局収支はどうかにはまだ験算が必要。}。七対子ドラドラ・字牌待ちなら立直\multicitep{miinin-stat, pp.~60--62; miinin-choice, pp.~103--104}、七対子のみ・待ち牌は端牌ではなくても立直\citep[][pp. 101--102]{miinin-choice}。

\subsubsection*{1枚切れの字牌にしろ}

\textbf{中張牌のドラより1枚見えの字牌待ちの方がいい}\citep[][pp. 109--110]{miinin-choice}。\textbf{タンヤオが見えそうにも関わらず1枚切れの字牌単騎にする}\citep[][pp. 111-112]{miinin-choice}。よって、手牌\ref{mj:rch-chitoi-1}と\ref{mj:rch-chitoi-2}は \mj{7s} 切り立直。

\tehai{3388m 11p 22337s 36 -6z}[東1局　南家　8巡目　ドラ \mj{7s}][rch-chitoi-1]

\tehai{3388m 2266p 2237s 7z -3s}[東1局　南家　8巡目　ドラ \mj{1z}][rch-chitoi-2]

\subsubsection*{字牌がドラの場合}

\textbf{字牌単騎にするとき字牌がドラでも立直する。ドラの字牌とドラではない字牌から選ぶ時でもドラ単騎にする}\citep[][pp.~113--114]{miinin-choice}。ドラ字牌のアガリ率（35\%）は非ドラ字牌（57\%）より低くても、表~\ref{tab:ev-chitoi} のように局収支から考えたらドラ字牌単騎の方がいい。よって、手牌\ref{mj:rch-chitoi-3}と\ref{mj:rch-chitoi-4}どちらもドラ \mj{2z} 単騎にして立直する。非ドラ字牌が1枚切れでも結論が変わらない。

\tehai{1133m 166p 44556s 2z -1p}[東1局　南家　8巡目　ドラ \mj{2z}][rch-chitoi-3]

\tehai{1133m 166p 4455s 23z -1p}[東1局　南家　8巡目　ドラ \mj{2z}][rch-chitoi-4]

\begin{table}[ht]
  \centering
  \begin{tabular}{l|ll}
    局収支 & 生牌 & 1枚切れ \\
    \hline
    非ドラ & 1700 & 2200 \\
    ドラ　 & 3000 &  \\
  \end{tabular}
  \caption{七対子字牌単騎の局収支}
  \label{tab:ev-chitoi}
\end{table}

\subsubsection*{中張牌ドラと生牌字牌}

中張牌ドラ七対子の局収支は1400\citep[][p.~110]{miinin-choice}。表~\ref{tab:ev-chitoi} によると非ドラ生牌字牌（1700）より低いので生牌字牌にする\footnote{中張牌ではないドラと生牌字牌の比較なら？これは作者が実戦で出会ったケース}。よって、手牌\ref{mj:rch-chitoi-5}は \mj{7s} を切って立直する\footnote{これも作者が実戦で出会ったケース。}。

\tehai{1105m 11p 11557s 334z}[東1局　南家　7巡目　ドラ \mj{7s}][rch-chitoi-5]

\subsubsection*{筋と字牌の選択}

\textbf{要すると「筋19 $\approx$ 1枚切れ字牌 ＞ 筋28 $\approx$ 生牌字牌 ＞ 筋37」}\citep[][pp.~115--116]{miinin-choice}。

%https://www.bilibili.com/read/cv15365496/?spm_id_from=333.999.0.0


\end{document}